\documentclass[a4paper,12pt,UTF8,oneside,hyperref]{ctexart} %twoside可设置双面打印格式;
\usepackage{times}%全文英文字体使用Times New Roman；
\usepackage[english]{babel} %设置英文单词的断行规则；
\linespread{1.5}%全文1.5倍行距；
\newfontface\englishbf{Times New Roman}[Contextuals={WordInitial,WordFinal}] %设置英文粗体方式；
%========================电子文档与超链接======================
\usepackage[backref=page,hyperfootnotes=false]{hyperref}
\usepackage[a4paper,left=2.5cm,right=2.5cm,top=2.5cm,bottom=2.5cm]{geometry}
\usepackage{graphicx}
\hypersetup{
 bookmarks = true,
 bookmarksopen = true,
 pdfborder=0 0 0,
 hyperfootnotes=false,
 pdfpagemode = FullScreen,
 pdfstartview = Fit,
 pdftitle = {coursework1_report_CID},
 pdfauthor = {Bowen Zhi}
}
\pdfbookmark[1]{Contents}{contents}
%===============================================================

%=========================设=置=页=眉===========================
\usepackage{fancyhdr} %为了保证正确长度的页眉下划线,fancyhdr宏包要在geometry宏包后使用；
\pagestyle{fancy}
\fancyhf{}
\lhead{\includegraphics[scale=0.04]{./pictures/ICL_logo.png}}
\usepackage{lastpage}
\rhead{\small\englishbf Reinforcement Learning Coursework(Report)}
%===============================================================

%==========================章节格式=============================
\usepackage{titlesec}
\usepackage{indentfirst}
\setlength{\parindent}{2em}
\titleformat{\section}[hang]{\englishbf\Large}{Chapter\ \thesection \quad}{0pt}{}{}
\titleformat{\subsection}[hang]{\englishbf\normalsize\raggedright}{\thesubsection\quad}{0pt}{}{}
\titleformat{\subsubsection}[hang]{\normalsize\kaishu\raggedright}{\thesubsubsection\quad}{0pt}{}{}
\titlespacing*{\section}{0pt}{0pt}{19.2pt}[0pt]
\titlespacing*{\subsection}{0pt}{0pt}{0pt}[0pt]
\titlespacing*{\subsubsection}{0pt}{0pt}{0pt}[0pt]
%===============================================================

%============================目录===============================
\usepackage{titletoc}
\titlecontents{section}[1.2em]{}{}{}{}[]  % 修改section的缩进
\dottedcontents{section}[1.2em]{\normalsize}{1.0em}{4pt}  % 修改section的起始缩进
\dottedcontents{subsection}[4.0em]{\normalsize}{2.0em}{4pt}
\dottedcontents{subsubsection}[7.0em]{\normalsize}{3.0em}{4pt}
%===============================================================

%==========================脚注格式=============================
\usepackage{pifont}
\renewcommand\thefootnote{\ding{\numexpr171+\value{footnote}}}% 更好看的一个脚注格式
\usepackage[perpage]{footmisc}%脚注每页清零编号
%===============================================================

%==========================插入图片=============================
\usepackage{graphicx}
% \usepackage{subfigure}%两图并排
\usepackage{picinpar}%文字绕排
\usepackage{caption}
\graphicspath{{./pictures/}}%设置图片放置路径
\captionsetup{font=small,labelfont=bf,labelsep=quad}
\captionsetup{skip=0pt}%标题与浮动环境内容的垂直间距无额外间距（设置为0pt）;
%===============================================================

%==========================插入表格=============================
\usepackage{booktabs,multirow} %三线表
\usepackage{longtable}         %长表格
\usepackage{bigstrut}
\usepackage{subcaption}        %插入并排图片
\usepackage{diagbox}
\usepackage{makecell}
\usepackage{tabularx}
\newcolumntype{Y}{>{\centering\arraybackslash}X}
%==========================解决三线表竖线过短的问题=============================
\renewcommand{\cmidrulesep}{0mm}  %定义两条相邻\cmidrule之间的间隔
\setlength{\aboverulesep}{0mm}    %在线条[不包括\toprule]上面增加一段垂直距离，此处为0mm
\setlength{\belowrulesep}{0mm}    %在线条[不包括\bottomrule]下面增加一条垂直距离，此处为0mm
\setlength{\abovetopsep}{0cm}     %在线条\toprule上面，即表格与上面的文字之间的距离。
\setlength{\belowbottomsep}{0cm}  %在线条\bottomrule下面，即表格与下面的文字之间的距离。
%===============================================================

%==========================数学公式=============================
\usepackage{amsmath}%公式
\usepackage{amssymb}
\usepackage{bm}%粗体
\usepackage{upgreek}%直立希腊字母
\usepackage{mhchem}%化学公式 \ce{SiO2} 或 \ce{H2O}
\DeclareMathOperator\dif{d\!}%微分算子\frac{\dif y}{dif x}
%\usepackage{xfrac}%行内小分式
\newenvironment{sequation}{\begin{equation}\small}{\end{equation}}
\newenvironment{tequation}{\begin{equation}\tiny}{\end{equation}}
%===============================================================

%==========================使用彩色=============================
%\usepackage[usenames,dvipsnames]{color}
\usepackage{color}
\usepackage[table]{xcolor}
\usepackage{colortbl}
\definecolor{mygray}{gray}{.9}
\definecolor{mypink}{rgb}{.99,.91,.95}
\definecolor{mycyan}{cmyk}{.3,0,0,0}
%================================================================

%==========================程序代码=============================
\usepackage{listings}
\usepackage{pythonhighlight}
%===============================================================

%=============================其它===============================
\usepackage{siunitx}%摄氏度 \SI{100}{\degreeCelsius}
\makeatletter%罗马数字
\newcommand{\rmnum}[1]{\romannumeral #1}% 罗马数字
\newcommand{\Rmnum}[1]{\expandafter\@slowromancap\romannumeral #1@}%罗马数字
\makeatother%罗马数字 \Rmnum{1} 或 \Rmnum{2} 或 \rmnum{3}
\usepackage{xfrac}%行内小分式
%==================================================================

%=============================文献格式===============================
\usepackage[numbers,authoryear]{natbib}         % 引用哈佛格式
\renewcommand{\bibfont}{\small\rmfamily} 
\setlength\bibsep{0pt}                          % 取消不同文献条目之间的距离
\usepackage{url}
%==================================================================

%=========================一个更好的下划线设置=======================
\usepackage{ulem}
\usepackage{enumitem}
\usepackage{hhline}% 双线表
\newcommand\dlmu[2][4cm]{\hskip1pt\underline{\hb@xt@ #1{\hss#2\hss}}\hskip3pt}
%==================================================================

%=========================公式按章编号===============================
\makeatletter
\@addtoreset{equation}{section}
\makeatother
\renewcommand{\theequation}{\arabic{section}-\arabic{equation}}

\makeatletter
\@addtoreset{figure}{section}
\makeatother
\renewcommand{\thefigure}{\thesection\,-\,\arabic{figure}}

\makeatletter
\@addtoreset{table}{section}
\makeatother
\renewcommand{\thetable}{\thesection\,-\,\arabic{table}}
%==================================================================

\usepackage{ifpdf}

\newtheorem{thm}{定理}

%==============积分符号修改，改原倾斜的积分符号为直立体==================
\usepackage{amsmath,amssymb}
\DeclareSymbolFont{EulerExtension}{U}{euex}{m}{n}
\DeclareMathSymbol{\euintop}{\mathop} {EulerExtension}{"52}
\DeclareMathSymbol{\euointop}{\mathop} {EulerExtension}{"48}
\let\intop\euintop
\let\ointop\euointop
%==================================================================

\usepackage{doc}
\begin{document}
\begin{titlepage}
  \begin{center}
  
  % 帝国理工校徽
  \includegraphics[scale=0.07]{./pictures/ICL_arms.png}\\[1cm]
  
  % 校名
  \LARGE \englishbf{Imperial College London} \\[0.5cm]
  
  % 系所信息
  \Large \englishbf{Department of Bioengineering} \\[1.5cm]
  
  % 标题
  \LARGE \englishbf{Reinforcement Learning Coursework (Report)} \\[1.5cm]
  
  % 个人信息表
  \renewcommand{\arraystretch}{1.2} % 行距
  \begin{table}[htbp]
  \centering
  \Large
  \begin{tabular}{c}
      \englishbf{Bowen Zhi} \\ 
      \englishbf{CID: } \\ 
      \englishbf{MSc in Human and Biological Robotics} \\
  \end{tabular}
  \end{table}
  
  \vfill
  
  % 日期
  \Large {November 2024}
  
  \end{center}
  \end{titlepage}
  
\clearpage
\phantom{s}
%\thispagestyle{empty}
%\newpage

\setcounter{page}{1}
\pagenumbering{Roman}
\cfoot{\small\thepage}

\vspace*{-21.6pt}
\begin{center}
\Large\englishbf
\rmfamily Template for the Thesis of Imperial College London
\end{center}
\par
\phantomsection
\addcontentsline{toc}{section}{ABSTRACT}
\begin{center}
    \fontsize{14pt}{17pt}\selectfont\rmfamily\bfseries
ABSTRACT
\end{center}

\par
This is abstract. The appendix of this article is used for example, like this "See Appendix \ref{matlab} for details".

\bigskip
\noindent\textbf{Key Words:}\quad Key1\quad Key2\quad Key3\quad Key4

\newpage

\phantomsection %设置“幻影标题”，解决额外插入目录（\addcontentsline）后总是引用到第一页的问题；
\addcontentsline{toc}{section}{Contents}
\vspace*{-21.6pt}
\tableofcontents
\newpage

\setcounter{page}{1}
\pagenumbering{arabic}
\cfoot{\small Page\, \thepage\,of\, \pageref{LastPage}}
\newpage\vspace*{-21.6pt}

\section{Introduction}
The format of this template is for reference only, the content is only for the purpose of displaying the format and is meaningless. For examlpe, we can have a matlab example in \ref{matlab}.

\newpage\vspace*{-21.6pt}
\section{Dynamic Programming}
\subsection{Parameterisation}

\cite{kaelbling1996reinforcement} discusses central issues of reinforcement learning, including trading off exploration and exploitation, establishing the foundations of the field via Markov decision theory, learning from delayed reinforcement, constructing empirical models to accelerate learning, making use of generalization and hierarchy, and coping with hidden state. Python example in \ref{python}。

They \citep{kaelbling1996reinforcement} discusses central issues of reinforcement learning, including trading off exploration and exploitation, establishing the foundations of the field via Markov decision theory, learning from delayed reinforcement, constructing empirical models to accelerate learning, making use of generalization and hierarchy, and coping with hidden state.

\subsection{Graphical representation}

The Imperial College's Coat of arms is shown in Figure \ref{Arm}.

\begin{figure}[!h]
	\centering
	\includegraphics[width=.35\textwidth]{ICL_arms.png}
	\caption{ICL arms}
	\label{Arm}
\end{figure}

\newpage

\phantomsection
\addcontentsline{toc}{section}{Reference}
\nocite{*}
\bibliographystyle{dcu}  % 使用Harvard格式
\bibliography{./reference/thesis}  % 指向bib文件

\titleformat{\section}[hang]{\englishbf\large}{Appendix\,\thesection \quad}{0pt}{}{}
\appendix
\newpage
\vspace*{-21.6pt}
	
\begin{center}
	\Large\englishbf
	\phantomsection
	\addcontentsline{toc}{section}{Appendix}
	Appendix
\end{center}

\section{Matlab example}
\label{matlab}
\definecolor{DarkGreen}{rgb}{0.0,0.4,0.0}
\lstloadlanguages{Matlab}
\lstset{language=Matlab,
	frame=single,                           
	rulesepcolor= \color{gray},
	breaklines=true
	basicstyle=\small\ttfamily,
	keywordstyle=[1]\color{blue}\bfseries,  % primitive funs in bold blue
	keywordstyle=[2]\color{purple},         % args of funs in purple
	keywordstyle=[3]\color{blue}\underbar,  % user funs in blue with underbar
	stringstyle=\color{purple},             % strings in purple
	showstringspaces=false,
	identifierstyle=,
	commentstyle=\usefont{T1}{pcr}{m}{sl}\color{DarkGreen}\small,
	tabsize=4,
	% more standard MATLAB funcs
	morekeywords={sawtooth, square},
	% args of funcs
	morekeywords=[2]{on, off, interp},
	% user funcs
	morekeywords=[3]{FindESS, homework\_example},
	morecomment=[l][\color{blue}]{...},     % line continuation (...) like blue comment
	numbers=left,
	numberstyle=\tiny\color{blue},
	firstnumber=1,
	stepnumber=1,
	escapeinside=``,
}

\begin{lstlisting}[language=Matlab]
	function f=det(P)
	n=size(P,4);
	f=0;
	if n==2  
	f=conv(P(:,:,1,1),P(:,:,2,2))-conv(P(:,:,1,2),P(:,:,2,1));   
	end
\end{lstlisting}


\section{Python example}
\label{python}
\begin{python}
	import math
	def hix(x):
	if x > 0:
	return x
	else:
	return 0
	
	def trans1(x1, x2, n, s1, s2, lock):
	# 1.	(A1,B1) -> (A1,B1)
	if lock == 1:
	dt = 22 * (math.floor(x1 / 2) + math.floor(x2 / 2))
	dw = math.floor(x1 / 2) * n + math.floor(x2 / 2) * (150 - n)
	return dt, dw, lock
	else:
	return 0, 0, lock
\end{python}
%\end{appendices}

\end{document}
